\documentclass[11pt]{article}
\usepackage[a4paper,margin=1in]{geometry}
\usepackage{subfiles}
\usepackage{amsmath,amssymb,bm}
\usepackage{hyperref}
\usepackage{booktabs}
\usepackage{array}
\usepackage{mathtools}

\title{Appendix 8 (to main theory): Frame Dragging (Lense--Thirring Effect) in One--Metric ISPG}
\author{}
\date{}

\begin{document}
\let\phi\varphi
\maketitle

\section*{Purpose and scope}
This appendix records the gravitomagnetic (frame-dragging) prediction of the main theory one-metric formulation in the slow-rotation, weak-field regime. The goal is a compact, citable statement of the Lense--Thirring precession rate and its dependence on the PPN/gravitomagnetic sector. At leading (1.5PN) order the result matches GR.

\section*{Non-negotiable consistency with the main theory (one-metric formulation)}
We work strictly within the main theory posture:
\begin{itemize}
\item matter is minimally coupled to the single physical metric $g_{\mu\nu}$ (no disformal mapping and no frame change);
\item the weak-field sector satisfies the standard PPN relations $\gamma=\beta=1$ (Appendix~4);
\item the distinguishing ISPG signatures are pushed to higher order (2PN), while the 1PN/1.5PN gravitomagnetic sector remains GR-identical under minimal coupling.
\end{itemize}
This scope covers Secs.~\ref{sec:metric_app8}--\ref{sec:LT} and the Earth-satellite numerical check in Sec.~\ref{sec:num}. The rotational-sector material in Sec.~\ref{sec:mondrot} is flagged explicitly below as a speculative MOND-bridge \emph{extension beyond} the declared MAIN backbone action and is not used in the 1.5PN GR-matching result.

\section{Weak-field slow-rotation metric and the gravitomagnetic potential}\label{sec:metric_app8}
In the weak-field, slow-motion expansion, a stationary rotating source generates an off-diagonal metric component $g_{0i}$ (gravitomagnetism). In standard PPN notation, the metric can be written as
\begin{align}
ds^2 &= -\left(1-\frac{2U}{c^2}\right)c^2dt^2 - \frac{4}{c^3}\,(\bm{A}\cdot d\bm{x})\,c\,dt
+ \left(1+\frac{2\gamma U}{c^2}\right)d\bm{x}^2 + O(c^{-4}),
\label{eq:ppn_gm}
\end{align}
where $U(\bm{x})$ is the Newtonian potential and $\bm{A}(\bm{x})$ is the gravitomagnetic vector potential. Appendix~4 gives $\gamma=1$ in the main theory formulation, so the spatial coefficient matches GR.

For an isolated spinning body of angular momentum $\bm{J}$, the external gravitomagnetic potential at leading order is
\begin{equation}
\bm{A}(\bm{x}) = \frac{G}{2}\,\frac{\bm{J}\times \bm{x}}{r^3},
\qquad r\equiv |\bm{x}|,
\label{eq:A}
\end{equation}
which implies the gravitomagnetic field
\begin{equation}
\bm{B}_g \equiv \nabla\times \bm{A} = \frac{G}{2r^3}\left[3(\bm{J}\cdot \hat{\bm{r}})\hat{\bm{r}}-\bm{J}\right].
\label{eq:Bg}
\end{equation}

\paragraph{Why this matches the main theory.}
In the one-metric minimal-coupling formulation, the $g_{0i}$ sector at 1.5PN order is taken to be governed by the same momentum-constraint structure as in GR, and no preferred-frame parameters are introduced. Together with $\gamma=1$, this yields the standard external solution \eqref{eq:A} at the order relevant to Lense--Thirring tests. (A full preferred-frame/vector-PPN derivation showing $\alpha_1=\alpha_2=\alpha_3=0$ for this theory would tighten the match; here we adopt the standard 1.5PN GR-sector assumption and verify that the resulting Lense--Thirring formula is consistent with current Gravity~Probe~B/LAGEOS data.)

\section{Lense--Thirring precession rates}\label{sec:LT}
\subsection{Gyroscope (spin) precession}
A gyroscope moving with velocity $\bm{v}$ in a stationary gravitomagnetic field experiences a frame-dragging precession
\begin{equation}
\bm{\Omega}_{\rm LT} = \frac{2}{c^2}\,\bm{B}_g
= \frac{G}{c^2 r^3}\left[3(\bm{J}\cdot \hat{\bm{r}})\hat{\bm{r}}-\bm{J}\right],
\label{eq:OmegaGyro}
\end{equation}
which is the standard Lense--Thirring result at leading order.

\subsection{Orbital node precession}
For a test body on an orbit of semi-major axis $a$ and eccentricity $e$ around a spinning mass, the secular precession of the ascending node is
\begin{equation}
\dot{\Omega}_{\rm node} = \frac{2GJ}{c^2 a^3(1-e^2)^{3/2}},
\label{eq:node}
\end{equation}
and the argument of pericenter receives the companion gravitomagnetic precession
\begin{equation}
\dot{\omega}_{\rm LT} = -\frac{6GJ\cos i}{c^2 a^3(1-e^2)^{3/2}},
\label{eq:pericenterLT}
\end{equation}
where $i$ is the inclination of the orbital plane relative to $\bm{J}$.

\section{Rotational-sector bridge to the MOND appendix (speculative extension; not part of the verified backbone)}\label{sec:mondrot}

\paragraph{Epistemic status of this section.}
The equations introduced in this section (Eqs.~\eqref{eq:Xdef_rot}--\eqref{eq:mu_rot}) constitute a \emph{speculative constitutive parametrization} for a rotational-sector MOND bridge. They are \emph{not} part of the MAIN backbone action declared in MAIN Sec.~2.1, which is the minimal $R+(\partial\phi)^2$ scalar-tensor form with minimal matter coupling. In particular:
\begin{itemize}
\item Eq.~\eqref{eq:Srot_eff} introduces a new non-minimal coupling between the local rotational kinetic invariant $X_{\rm rot}$ (defined in Eq.~\eqref{eq:Xdef_rot}; opposite sign to the main Horndeski $X$) and the gravitomagnetic invariant $|\bm{B}_g|^2$. This is a term \emph{beyond} the declared backbone; its status as a consistent EFT extension (ghost freedom, energy conditions, classification within Horndeski / beyond-Horndeski) is \emph{not} verified here and is deferred to future work.
\item Eq.~\eqref{eq:g_split_rot} is a \emph{postulate} on how the total acceleration $g$ decomposes into a Newtonian channel $g_N$ and a halo-effective channel $g_h$; it is not derived from Eq.~\eqref{eq:A_mond_rot}. The interpolating form $\mu(x)=x/(1+x)$ in Eq.~\eqref{eq:mu_rot} therefore follows from this postulate in one step, and is not an independent derivation from the backbone.
\item The 1.5PN GR-matching result of Sec.~\ref{sec:LT} does \emph{not} depend on this bridge: setting $Z_{\rm rot}\to 0$ in the strong-field limit (below) leaves the assumed 1.5PN GR sector intact. The Solar-System tests are controlled exclusively by Secs.~\ref{sec:metric_app8}--\ref{sec:LT}.
\item The companion MOND paper~\cite{Lotishvili2026MOND} relies on this parametrization. The cross-paper ``builds strictly on MAIN backbone'' statement should therefore be read with this qualification: the MOND sector requires an explicit rotational-operator extension beyond the MAIN minimal action, whose full derivation from first principles is an open task.
\end{itemize}

With this status recorded, the parametrization proceeds as follows. We use a local positive-sign kinetic invariant $X_{\rm rot}$ (the standard main-text/Appendix~11 Horndeski convention is $X = -\tfrac{1}{2}g^{\mu\nu}\partial_\mu\varphi\,\partial_\nu\varphi$; the local rotational-bridge definition below has opposite sign for notational convenience and does not enter the diagonal scalar sector):
\begin{equation}
X_{\rm rot} \equiv \frac12 g^{\mu\nu}\partial_\mu\phi\,\partial_\nu\phi,
\qquad
X_0 \equiv \frac12\left(\frac{2a_0}{c^2}\right)^2,
\label{eq:Xdef_rot}
\end{equation}
so that in the weak field $\sqrt{X_{\rm rot}/X_0}=g/a_0$. Introduce the rotational stiffness
\begin{equation}
K_{\rm rot}\!\left(\frac{X_{\rm rot}}{X_0}\right) \equiv 1+\sqrt{\frac{X_{\rm rot}}{X_0}}
= 1+\frac{g}{a_0},
\qquad
Z_{\rm rot} \equiv K_{\rm rot}^{-1}.
\label{eq:Krot}
\end{equation}
The interpretation is simple: $K_{\rm rot}$ measures how ``loaded'' or sticky the medium is against twisting, whereas $Z_{\rm rot}$ is the rotationally free fraction. Strong static deficit means large $K_{\rm rot}$ and small $Z_{\rm rot}$; weak static deficit means $K_{\rm rot}\to 1$ and $Z_{\rm rot}\to 1$.

A minimal effective completion that preserves the diagonal scalar sector is then
\begin{equation}
S_{\rm rot}^{\rm eff}
= \frac{1}{16\pi G}\int d^4x\,\sqrt{-g}\,
Z_{\rm rot}\!\left(\frac{X_{\rm rot}}{X_0}\right)\,\mathcal I_{\rm rot},
\label{eq:Srot_eff}
\end{equation}
where $\mathcal I_{\rm rot}$ is a positive invariant built only from the rotational sector. The Hubble scale is already encoded through $a_0=cH/(2\pi)$ inside $X_0$. In the weak-field stationary limit one may take
\begin{equation}
\mathcal I_{\rm rot}\propto \frac{|\nabla\times \bm{A}|^2}{c^2}
= \frac{|\bm{B}_g|^2}{c^2}.
\label{eq:Irot_wf}
\end{equation}

Varying the extra term with respect to $\bm{A}$ gives
\begin{equation}
\delta S_{\rm rot}^{\rm eff}
\propto
\int d^3x\,
\bigl[\nabla\times(Z_{\rm rot}\bm{B}_g)\bigr]\cdot \delta\bm{A},
\end{equation}
so the stationary weak-field rotational equation takes the schematic form
\begin{equation}
\mathcal D_{\rm GR}[\bm{A}]
\;+\;
\nabla\times(Z_{\rm rot}\bm{B}_g)
\;=\;
\frac{16\pi G}{c^2}\,\bm{J},
\label{eq:A_mond_rot}
\end{equation}
where $\mathcal D_{\rm GR}[\bm{A}]$ denotes the usual GR operator from the Einstein--Hilbert sector and $\bm{J}$ is the angular-momentum current density. In Coulomb gauge and for slowly varying $Z_{\rm rot}$, Eq.~\eqref{eq:A_mond_rot} reduces approximately to a GR vector equation plus an additive term proportional to $-Z_{\rm rot}\nabla^2\bm{A}$.

\paragraph{Immediate consequence.}
Because $0\leq Z_{\rm rot}\leq 1$, the operator-level enhancement of the GR-like rotational response is only $O(1)$ (at most a factor of order two in the slowly varying limit). Therefore this extension, by itself, does not yet bridge the old perturbative amplitude gap if the right-hand-side source remains the ordinary GR mass-current source. A full MOND completion likely requires not only operator renormalization but also a new constitutive \emph{source-side} response of the rotating medium.

Because $Z_{\rm rot}$ depends on $X_{\rm rot}$, the same extension also feeds back into the scalar equation:
\begin{equation}
\Box\phi
\;=\;
\text{(matter source)}
\;+\;
\frac{1}{X_0}\,
\nabla\!\cdot\!\left[
\mathcal I_{\rm rot}\,
\frac{dZ_{\rm rot}}{d(X_{\rm rot}/X_0)}\,
\nabla\phi
\right],
\label{eq:phi_backreaction_rot}
\end{equation}
and since $dZ_{\rm rot}/d(X_{\rm rot}/X_0)<0$, the rotational sector deepens rather than erases the low-pressure state in the MOND regime.

\paragraph{Postulate (not derivation).}
We \emph{postulate} that the total inward support decomposes into a Newtonian channel $g_N$ and a halo-effective channel $g_h$ weighted by $Z_{\rm rot}$:
\begin{equation}
g_N = (1-Z_{\rm rot})\,g,
\qquad
g_h = Z_{\rm rot}\,g,
\label{eq:g_split_rot}
\end{equation}
which immediately yields
\begin{equation}
\mu(x)=\frac{g_N}{g}=\frac{x}{1+x},
\qquad x\equiv \frac{g}{a_0}.
\label{eq:mu_rot}
\end{equation}
Eq.~\eqref{eq:g_split_rot} is \emph{not} derived from Eq.~\eqref{eq:A_mond_rot}; it is an independent constitutive postulate that, together with Eq.~\eqref{eq:Krot}, fixes the $\mu(x)=x/(1+x)$ interpolation in one step. The companion MOND paper~\cite{Lotishvili2026MOND} reobtains the same local ratio from a bound-structure source-side vortex closure; that derivation is the primary locus where $\mu(x)$ is fixed, and the present appendix records Eqs.~\eqref{eq:g_split_rot}--\eqref{eq:mu_rot} only as the compact rotational-slot parametrization consistent with that closure.

\paragraph{Strong-field limit and Solar-System numerical check.}
For $g\gg a_0$, one has $K_{\rm rot}\simeq g/a_0\gg 1$ and $Z_{\rm rot}\simeq a_0/g\ll 1$, so the \emph{additional} rotational-MOND term \eqref{eq:Srot_eff} switches off and the usual Lense--Thirring sector from the Einstein--Hilbert part is recovered.

Concretely, with $a_0\simeq 1.2\times 10^{-10}\,{\rm m\,s^{-2}}$, the Sun-at-Earth Newtonian acceleration is $g_\odot\simeq 6.0\times 10^{-3}\,{\rm m\,s^{-2}}$, giving
\begin{equation}
\frac{g_\odot}{a_0}\simeq 5.0\times 10^{7},
\qquad
Z_{\rm rot}\simeq \frac{a_0}{g_\odot}\simeq 2.0\times 10^{-8},
\end{equation}
so the rotational-MOND correction to the Lense--Thirring rate at Earth-orbit distance from the Sun is bounded by $\sim 2\times 10^{-8}$ of the GR value --- well below current observational sensitivity. For a LAGEOS-like satellite in Earth's field, $g_\oplus/a_0\sim 10^{11}$, giving $Z_{\rm rot}\sim 10^{-11}$. The extension is therefore numerically inert for all Solar-System frame-dragging tests, and the 1.5PN GR-matching statement of Sec.~\ref{sec:LT} is preserved at the relevant observational precision.

\paragraph{Status of this bridge.}
Equations~\eqref{eq:Krot}--\eqref{eq:mu_rot} should be read as the compact rotational-sector parametrization passed forward to the companion MOND paper~\cite{Lotishvili2026MOND}, not as a verified consequence of the MAIN backbone action. Their full derivation from first principles --- including the classification of \eqref{eq:Srot_eff} as a consistent EFT extension, the ghost analysis, the energy-condition analysis, and a first-principles derivation of the constitutive split \eqref{eq:g_split_rot} --- is an open task. The present appendix records Sec.~\ref{sec:mondrot} only to make the rotational slot explicit and to confirm that its $Z_{\rm rot}\to 0$ limit in the Solar System leaves the recovered 1.5PN gravitomagnetic result (under the stated GR-sector assumption) intact.

\section{Numerical scale (Earth satellites)}\label{sec:num}
For Earth, $J_\oplus\simeq 5.86\times 10^{33}\,{\rm kg\,m^2/s}$. For a LAGEOS-like orbit ($a\simeq 1.227\times 10^7\,{\rm m}$, $e\simeq 0.004$), Eq.~\eqref{eq:node} gives a nodal precession of order
\begin{equation}
\dot{\Omega}_{\rm node}\sim 30~{\rm mas/yr},
\end{equation}
consistent with the classic GR target of Lense--Thirring satellite tests.

\section*{Takeaway}
In the main theory one-metric ISPG formulation, the weak-field gravitomagnetic sector at 1.5PN order is recovered under the stated 1.5PN GR-sector assumption (Secs.~\ref{sec:metric_app8}--\ref{sec:LT}, numerical check in Sec.~\ref{sec:num}; a full preferred-frame/vector-PPN derivation showing $\alpha_1=\alpha_2=\alpha_3=0$ would tighten this match). The predicted frame-dragging effects at this order are the standard Lense--Thirring precessions \eqref{eq:OmegaGyro}--\eqref{eq:pericenterLT}.

Section~\ref{sec:mondrot} records a rotational-sector MOND-bridge parametrization that is \emph{not} part of the MAIN backbone action: it introduces a new scalar--gravitomagnetic coupling \eqref{eq:Srot_eff} and a constitutive split postulate \eqref{eq:g_split_rot} whose first-principles derivation and EFT-consistency analysis are deferred to future work. The bridge is numerically inert in the Solar System ($Z_{\rm rot}\lesssim 10^{-8}$), so its status does not affect the recovered GR-matching result (under the stated assumption). Distinctive ISPG discriminators beyond the assumed-GR-limit regime therefore reside either in the rotational MOND sector (pending full derivation) or at higher (2PN) order (pending the open-task coefficient of Appendix~6).

\end{document}
